\chapter{Static navigation analysis and scope attribution}
\label{ch:analysis}

This chapter presents the analysis that is the main contribution of the thesis.
Section~\ref{sec:model} defines the graph model and states the soundness property
claimed for it. Section~\ref{sec:attribution} presents scope attribution, the analysis
proposed in answer to RQ3. Section~\ref{sec:implementation} describes how both are
implemented in the extractor. The evaluation of the analysis is described separately,
in Chapter~\ref{ch:evaluation}.

\section{Graph model and soundness}
\label{sec:model}

\begin{definition}[UI transition graph]
A UI transition graph is a labelled transition system $(S, \Sigma, G, M, E)$ with
\[
  E \;\subseteq\; S \times \Sigma \times G \times M \times S .
\]
$S$ is the finite set of \emph{screens}, the states of the system and the nodes of the
graph. $\Sigma$ is the finite set of \emph{event kinds} that can be raised at a screen,
such as the activation of a \texttt{routerLink} or a call to \texttt{Router.navigate}.
$G$ is the set of \emph{guards}, each held as the symbolic source text of the condition
under which an edge may be taken. $M$ is the set of \emph{modalities}, which has
exactly three elements,
\[
  M \;=\; \{\mathit{must}, \mathit{may}, \mathit{unknown}\},
\]
and whose elements are confidence labels: a modality states how firmly the extractor
asserts an edge, and never denotes a screen, a route or a navigation target. An edge
$(s, \sigma, g, m, s') \in E$ reads: at screen $s$, event $\sigma$ under guard $g$ leads
to screen $s'$, asserted with confidence $m$. Both the source $s$ and the target $s'$
are elements of $S$, namely the first and the last component of the tuple. Constructs
in the source that agree on source, event kind, guard and target yield the same edge, which
carries the strongest modality among them, so $E$ holds at most one modality for each
$(s, \sigma, g, s')$; the paragraph on witnesses below says how the constructs are traced.
\end{definition}

\textbf{What the materialised graph carries beyond the model.} The extractor emits two node
kinds that are not screens, and the edge relation above does not range over them. An
\emph{unplaced origin} stands for a component, with the navigation sites in it, that the
analysis could not place under any screen; a \emph{dynamic target} stands for a destination only the run time knows. Both are
kept in the output rather than dropped, because a site the analysis cannot place is a loss
Chapter~\ref{ch:results} has to count, and a destination it cannot name is one the run time
could still take. Table~\ref{tab:census} reports how much of each graph they are. Over the
corpus \CensusOtherNodes{} of the \CensusNodes{} nodes are of these two kinds,
\CensusOtherNodeShare{} of them, and \CensusTargetNotScreen{} of the \CensusEdges{} edges
reach a dynamic target. A further \CensusMaySourceNotScreen{} edges leave an unplaced origin
and are held at \emph{may}, so they sit inside a tier Chapter~\ref{ch:results} reports while
their source is not a screen. An edge that reaches a dynamic target asserts that the screen it
leaves has an outgoing navigation whose destination the analysis could not name, and asserts
nothing about which screen is reached. Claim~\ref{claim:must} and the conversion ratio on which RQ1 is
defined range over the edges whose source and target are both screens.

\begin{table}[t]\small
\caption{What the published graphs are made of: nodes by kind, where an anchor stands for a
whole subtree of screens and \emph{other} counts unplaced origins and dynamic targets
together, and edges by modality. These are the attributed graphs; the node column of
Table~\ref{tab:conversion} counts the graphs of the same snapshot without attribution, which
hold no anchors and fewer \emph{other} nodes.}
\label{tab:census}
\TabCensus
\end{table}

A \emph{must} edge is claimed to fire for every input on which the enclosing code
executes; a \emph{may} edge is possible but
conditional; \emph{unknown} marks an edge whose modality could not be determined. The
third modality carries the destinations the analysis cannot name. It labels only edges of the
materialised graph that lie outside $E$, namely edges into the
\texttt{u\_}\,sink of Section~\ref{sec:res-round2}, which records that a screen is left by
an event whose target only the run time knows. The pinned extractor does produce such edges:
\CensusUnknown{} of the \CensusEdges{} edges in the corpus, in \CensusUnknownApps{} of the
\CensusApps{} applications. Carrying this modality and reaching a node that is not a screen
coincide on every edge of every corpus graph, and the census refuses to report a graph
in which they come apart. The conversion ratio of RQ1 and the soundness property of
Claim~\ref{claim:must} are defined over the other two tiers, so the unknown tier is counted in
the edge totals and in no rate.
Guards are stored as source text and are never evaluated by the tool. A consumer of
the graph can read a guard, but is not expected to decide whether it holds. This keeps
the set $G$ opaque and means that the model does not need to represent application
state.

Alongside the modality the artefact records a numeric \emph{confidence} on each edge, which is
not part of the model and carries no claim this thesis makes. It is a finer version of the same
ordering: $1$ for a must edge; $0.6$ for a guarded edge with a literal target, or $0.5$ when
the guard is asynchronous; at most $0.5$ for every other may edge, whatever demoted it; and $0$ for an edge into a
dynamic target. Where several constructs yield one edge, the edge carries the modality and the
confidence of the most certain of them, and each construct it cites records the modality it
earns on its own, so the edge is as certain as the best evidence for it and no construct is cited
as stronger evidence than it is. Every figure in this thesis is reported over the modalities,
never over the confidence.

\textbf{Screens and routes.} The set $S$ is derived from the route configuration, and
not every route contributes a screen. Write $\mathit{Routes}$ for the set of routes the
application declares and $\mathit{component}(r)$ for the component that route $r$
renders. Then
\[
  S \;=\; \{\, \mathit{screen}(r) \;\mid\; r \in \mathit{Routes}
  \text{ and } \mathit{component}(r) \text{ is defined} \,\},
\]
where $\mathit{screen}(r)$ is the node for route $r$, identified by the full path that
concatenates the paths of $r$ and of its ancestors. A route that declares
\texttt{children} but no \texttt{component} is only a path prefix, as
Section~\ref{sec:angular} describes: it contributes no element of $S$, and its path is
carried into the full paths of its children. A route whose only effect is a
\texttt{redirectTo} renders nothing and likewise contributes no element of $S$. In Angular
a redirecting route is only a redirect: the router refuses a configuration that declares
\texttt{redirectTo} and \texttt{children} together, and likewise \texttt{redirectTo} with a
\texttt{component}, so a redirecting route can never also be the layout under which a child
renders. Both refusals were read off the router's own configuration validator in the version
vendored with this work, Angular~18.2.14; no corpus repository is checked out where the
invariant could be re-read, so the statement holds only for that version. The same is not true of every client-side router.
Vue Router applies a record's redirect only when that record is the last one matched, so a
redirecting parent does render its component for a matching child.
Section~\ref{sec:transfer} returns to this, because the model's step from
``redirects'' to ``renders nothing'' is a fact about Angular rather than about routers.
Because
$\mathit{screen}(r)$ is identified by the full path, the map $\mathit{screen}(\cdot)$ is
injective on the routes whose full paths are pairwise distinct. Angular also permits two
component-bearing routes with the same full path, told apart at run time by
\texttt{pathMatch} or by a \texttt{canMatch} guard. The model assumes that this does not
occur: the extractor keeps the route declared first under a full path and drops the
later ones. For Vue and for React the extractor also declares the drop, as a note carrying the
file and line of the record it discarded, so a graph built from a table with a repeated address
says so rather than being quietly one screen short; the note is true of the record and
over-stated about the graph, because where both records name the same component the address
still renders it. Angular's extractor drops silently, which is one of the asymmetries
Section~\ref{sec:transfer} returns to. Two routes with different paths that render the same component give two
screens, which is the case that makes the attribution of
Section~\ref{sec:attribution} necessary.

\textbf{Witnesses.} Every edge of static or runtime origin carries a witness: the file
and line of a construct that produced it, or an observation identifier. Where several
constructs produce the same edge, one of them is kept as the witness and the others are
recorded as merged witnesses, so that every construct behind an edge can be traced
(Section~\ref{sec:metrics}). No edge is
added to the base graph without a witness, and edges added manually are kept in a
separate overlay. The tool's specification states this as an invariant, and the first evaluation method checks it mechanically.

\begin{claim}[Soundness of the must tier]
\label{claim:must}
Every must edge corresponds to a navigation that occurs, for every input on which the
enclosing code executes, over the declared route set, under the route-order assumption of
Section~\ref{sec:approximations}, and for a site whose component is not placed by a join over an
ambiguous selector, within the following execution model: the event is a plain primary
activation that navigates the current browsing context, no other listener the event reaches
starts a navigation, and no call the reader cannot resolve to a body in the application, nor any
work the handler schedules for later, starts one.
\end{claim}

\noindent The execution model is part of the claim and not a detail of the implementation.
Section~\ref{sec:approximations} lists the cases it excludes: a call into a library that cancels
the navigation by starting another, a second listener on the element or an ancestor, and an
activation with a modifier key or into another browsing context. On such an execution a must edge
can name a navigation that does not complete, and the claim says nothing about it.

That assumption carries much of the tier: without it the must tier holds \RouteOrderMustWithout{}
of its \RouteOrderMust{} edges, on \RouteOrderSitesWithout{} of \RouteOrderSites{} constructs.

The evaluation tests a necessary condition of this claim rather than the claim itself,
and the distinction decides how the figures of Chapter~\ref{ch:results} may be read. The
witness check of Section~\ref{sec:protocol} re-reads the source at every location an edge
cites and asks whether the construct found there names the destination the edge names. A must
edge whose construct that earns must names a different destination cannot satisfy the claim,
so a failure of the check there refutes it; agreement does not establish it. The other
constructs a must edge cites are evidence for may and are not judged against the claim. Whether the navigation
occurs \emph{for every input on which the enclosing code executes} is decided by the rules
that place an edge in the must tier instead of the may tier. A site is demoted to may in the
following cases.
\begin{itemize}
\item \textbf{Guards and resolvers.} A guard the navigation runs demotes it, whatever the guard
would decide, since the guard half of the rules is read off the route configuration and never
evaluates a guard. Which guards a navigation runs is decided the way the router decides it. The
routes the source and the target share are reused unless their parameters can change or
\texttt{runGuardsAndResolvers} asks otherwise, and a reused route's \texttt{canActivate} and
resolvers do not run again. A navigation runs the \texttt{canMatch} of every route its address
is matched through; the \texttt{canActivate}, \texttt{canLoad} and resolvers of every route it
enters, \texttt{canLoad} until the route's code has been loaded, which the extractor, not knowing
whether it has, takes as every time; the \texttt{canActivateChild} of every route above those;
and the \texttt{canDeactivate} of every route it leaves. A child at the empty path that the
target's own address resolves to is one of the routes it enters. A resolver is awaited by the
router and demotes the edge as an \texttt{await} does, and a guard list the reader cannot bring
back to an array demotes it as a guard does.
\item \textbf{Target and source addresses.} The target is the address the source already
displays: the router ignores such a navigation in its default configuration, or follows it only
when a parameter or the query differs, so the extractor holds it at may as a matter of policy,
which is the safe side for an application that configures otherwise; a shell or anchor edge is
judged per screen and drawn per screen where the screens disagree. The source is never displayed
at its own address because a redirect answers it. The address is one that two declared routes,
or a route and a redirect, can answer, since Angular takes whichever it meets first. The base a
relative target resolves against cannot be read. The URL names an auxiliary outlet, whose route
the adapter does not read.
\item \textbf{Conditional execution.} The site runs only on some executions of its own handler:
under a template control-flow block, a structural directive, an \texttt{<ng-template>}, an ICU
case or a container that renders its content lazily; under a branch, in a loop or a
\texttt{catch} clause, in an optional call or a default value, inside a callback; after an
\texttt{await}, an early \texttt{return} or \texttt{throw} or a \texttt{break} out of its block;
in a function declared at module level or a member of its class that no handler calls
unconditionally; in a component mounted under a template condition on that screen; or in a route
component behind a conditional outlet. A handler is a member the framework calls itself, a
lifecycle hook or a member \texttt{@HostListener} registers, or one the template binds to an event
as a bare call outside any template condition.
\item \textbf{A later navigation may cancel it.} The router cancels a navigation when another
starts before it has finished, so a navigation is demoted when code that runs after it, before
the event is handled, may navigate. That code is the rest of the expression that holds it (the
body of a call it is an argument of, a right operand, a branch of a conditional whose condition it
is), the statements after it in its block and in every block around it, a \texttt{finally} block
around it, and what runs once its function returns: what each caller in the class does after
calling that function, handler or not, and for a constructor or lifecycle hook the hooks Angular
calls after it. Code there may navigate when it holds a navigation call, in a branch or a callback
as much as at the top, or a call to a member of the class, a function of the file or a method of
an injected service whose body does. An awaited navigation has finished before the rest of its
function runs, so what follows it there does not count, but its callers go on as soon as the
function first awaits, so what they do still counts. A call the reader cannot resolve to a body
in the application is taken not to navigate (Section~\ref{sec:approximations}).
\end{itemize}
An exception from an earlier synchronous
call is not counted, since every call can throw; an \texttt{await} is, because a rejected request
is an ordinary outcome of the code that awaits it. The router's own fetch of a lazily loaded
route is not counted either: a failed chunk load is a fault of the deployment rather than an
outcome of the application, and the extractor takes the code as deployed. The extractor's unit
tests cover each of these rules in both directions, a site the rule demotes and one it leaves
alone, and a mutation sweep deletes each rule in turn and checks that some test fails
(Section~\ref{sec:res-round2}); they share the extractor's reading of each construct, and no
second classifier re-derives them. What stands in for one, for the condition rules, is a review
in which the assistant the declaration describes read the constructs behind the round's must
tier in the source, twice and separately, without the extractor's readers, and recorded a verdict
for each (Section~\ref{sec:res-round2}). Section~\ref{sec:res-round2} reports the corrections to
the extractor's reading. Three made the tier smaller, each hidden by the one before; the fourth,
found by review rather than by the witness check, closed gaps in the demotion rules, replaced the
third's merge rule, which had demoted an edge one of its constructs made certain, and replaced a
guard rule that demoted an edge for guards the router does not run on it. What the
check establishes is therefore narrower than the claim: that every cited site exists in the
source and agrees on its target. The stronger half of the claim rests on the demotion rules
and is not independently verified.

This is the only soundness property claimed in the thesis, and it holds only relative
to the declared route set. No such property is claimed for the may and unknown tiers.
The witness of a may edge shows that the corresponding call site exists in the source,
but it does not show that the transition is ever taken. This difference between the
tiers matters for the evaluation. Whether the target found at a witness location agrees
with the target of the edge is a meaningful question for must edges, and for may edges
whose target is a literal. It is not a meaningful question for edges that the extractor
produced by over-approximation, and treating it as one leads to spurious failures, as
described in Section~\ref{sec:corrections}.

\section{Scope attribution}
\label{sec:attribution}

This section specifies more than the prototype contains, and the division is the same
throughout. Component nesting, router outlets, the shell rule, anchored edges and the call
graph with dependency injection, read at the granularity of class members with bindings taken
application-wide, are implemented in the extractor and are the parts evaluated in
Chapter~\ref{ch:evaluation}. The dispatch relations that carry NgRx effects and router event
subscriptions, guards that return a \texttt{UrlTree}, the injector hierarchy that decides which
binding of a token applies where, and navigation inside signal effects and observable
subscriptions, which the model labels \emph{unknown} (Section~\ref{sec:boundary}), are part of
the model only: they are specified here to define the analysis, and the extractor approximates
or omits them as Section~\ref{sec:approximations} lists. A navigation site no implemented
relation reaches contributes no edge, and the evaluation counts it in the defect class
``Attribution, service'' of Table~\ref{tab:loss-rules}.
Section~\ref{sec:implementation} gives the same division from the implementation side.

\subsection{The defect being addressed}

In the version of the extractor used for the baseline measurements, a navigation
site is attributed to a screen only if it occurs directly in the component that the
route for that screen renders. Navigation sites anywhere else are not attributed to any
screen. This one cause explains three symptoms in the baseline: shared components contribute no edges, because the application shell places the header next to the router outlet and not inside the subtree of any screen; a component mounted at several routes contributes one edge instead of one per route; and applications that navigate from injected services produce no edges at all.

\subsection{Scope}
\label{sec:scope}

Every navigation site is anchored at a \emph{scope}: the set of screens from which it
can fire. A scope is thus an element of the powerset $\mathcal{P}(S)$, and the abstract
domain of the analysis is $\mathcal{P}(S)$ ordered by inclusion, with set union as join,
\[
  X \sqsubseteq Y \iff X \subseteq Y, \qquad X \sqcup Y \;=\; X \cup Y ,
\]
which is the powerset lattice over the screens. Its least element is the empty scope,
\[
  \bot \;=\; \varnothing ,
\]
a site that fires from no screen at all. The greatest element is $\top = S$, every screen.

Two different things are called a scope, and the rest of this thesis keeps them apart.
Write $\mathit{scope}^{*}(x)$ for the \emph{true scope} of a site $x$, the set of screens
from which $x$ really can fire, which is a property of the application and not of any
analysis; and $\mathit{scope}(x)$ for the \emph{computed scope}, the set the fixpoint of
Section~\ref{sec:inputs} returns. A screen in
$\mathit{scope}^{*}(x) \setminus \mathit{scope}(x)$ is a scope too small and a screen in
$\mathit{scope}(x) \setminus \mathit{scope}^{*}(x)$ is a scope too large. No true scope is
established anywhere in this thesis. For one application a third set is read off the sources by
hand, the \emph{structural reference} $\mathit{scope}^{R}(x)$, which applies the analysis's
own relations without its implementation (Section~\ref{sec:res-groundtruth}). Agreement
between $\mathit{scope}$ and $\mathit{scope}^{R}$ shows that the relations are implemented as
stated; it does not show that either equals $\mathit{scope}^{*}$. Everywhere else the thesis
reports the computed scope and says so.

The empty computed scope is likewise two situations, and conflating them would hide a defect
inside a property of the application. A site is \emph{unreachable} when
$\mathit{scope}^{*}(x) = \varnothing$: the code exists and no screen can reach it, so
contributing no edge is the correct answer. A site is \emph{unresolved} when
$\mathit{scope}(x) = \varnothing$ because no rule of the analysis reaches it (a navigation
in a service nothing is known to inject, or behind a construct the reader could not follow),
and the true scope is simply unknown. Both yield $\bot$, so $\bot$ alone does not say which
holds. The extractor therefore does not let an unresolved site vanish silently. A site in a
component that no rule places is anchored at an unplaced origin, a node that carries the
component's file, and its edges are held at \emph{may}; a site in any other class whose scope
is empty contributes no edge, and the extractor records a note naming its file and the reason.
Either trace is what makes the unresolved sites countable in Chapter~\ref{ch:results} rather
than indistinguishable from code that is genuinely unreachable.

Three elements of the domain are named, because attribution anchors edges at them.
Writing $r' \succeq r$ for the reflexive ancestor relation on routes, that is $r'$ is
$r$ itself or an ancestor of $r$ in the route tree, they are
\[
  \{\mathit{screen}(r)\}, \qquad
  \mathit{subtree}(r) \;=\; \{\, \mathit{screen}(r'') \in S \;\mid\; r \succeq r'' \,\},
  \qquad \top \;=\; S ,
\]
and they form the chain
\[
  \bot \;\sqsubseteq\; \{\mathit{screen}(r)\} \;\sqsubseteq\; \mathit{subtree}(r)
  \;\sqsubseteq\; \top .
\]
The chain is stated for a route $r$ with $\mathit{component}(r)$ defined. For a prefix
route, which has \texttt{children} but no \texttt{component}, $\mathit{screen}(r)$ is
undefined and only $\mathit{subtree}(r)$, the set of the screens of its descendants, is
defined. The chain then reads $\bot \sqsubseteq \mathit{subtree}(r) \sqsubseteq \top$.
These names are a representation device for anchoring edges, not a restriction of the
domain. The join of two scopes is their union, not the least named element above both: a
component mounted at two unrelated leaf routes has the scope
$\{\mathit{screen}(r_1), \mathit{screen}(r_2)\}$, which is none of the three named
elements. The names matter for storage, where an edge anchored at $\top$ or at
$\mathit{subtree}(r)$ is kept once rather than once per screen.
Table~\ref{tab:scopes} assigns a scope to each Angular construct.
Figure~\ref{fig:scopes} shows the three kinds of scope in a small application, and
Figure~\ref{fig:lattice} shows the domain for three of its screens.

\begin{figure}[t]
\centering
\begin{tikzpicture}[node distance=9mm and 6mm]
  % shell
  \node[comp, minimum width=26mm] (app) {AppComponent};
  \node[comp, right=6mm of app, fill=accentlight, draw=accent] (header) {HeaderComponent};
  \draw[nestedge] (app) -- (header);
  \node[comp, below=6mm of app, fill=white, font=\scriptsize\ttfamily] (outlet) {<router-outlet>};
  \draw[nestedge] (app) -- (outlet);

  % screens
  \node[screen, below=11mm of outlet, xshift=-30mm, fill=warmlight] (home) {/};
  \node[screen, right=of home, fill=warmlight] (feed) {/feed};
  \node[screen, right=of feed] (settings) {/settings};
  \node[screen, below=11mm of settings] (profile) {/settings/profile};
  \foreach \s in {home, feed, settings} \draw[->, muted] (outlet.south) -- (\s.north);
  \draw[->, muted] (settings) -- node[lbl, right] {\scriptsize outlet} (profile);

  % shared card
  \node[comp, below=11mm of home, xshift=12mm, fill=warmlight, draw=warm] (card) {CardComponent};
  \draw[nestedge] (home) -- (card);
  \draw[nestedge] (feed) -- (card);

  % scopes
  \begin{scope}[on background layer]
    \node[scope, fit=(settings)(profile), inner sep=5pt, label={[font=\scriptsize, text=accent]right:$\mathit{subtree}(\text{/settings})$}] (sub) {};
    \node[draw=accent, rounded corners=6pt, densely dotted, fit=(home)(feed)(settings)(profile)(card)(sub), inner sep=7pt,
      label={[font=\scriptsize, text=accent, anchor=south west]north west:$\top$ = every screen}] (topbox) {};
  \end{scope}

  % legend notes
  \node[font=\scriptsize, align=left, anchor=west] at ($(header.east)+(4mm,0)$) {links here: scope $\top$,\\stored once};
  \node[font=\scriptsize, align=center, anchor=north, text=warm] at ($(card.south)+(0,-6mm)$) {links here: scope $\{\text{/}, \text{/feed}\}$,\\one edge per screen};
  \node[font=\scriptsize, align=left, anchor=west, text=accent] at ($(profile.east)+(3mm,-6mm)$) {links in the settings\\template: anchored at the\\subtree, stored once};
\end{tikzpicture}
\caption{Scopes in a small application. The header is rendered next to the outlet, so its
links fire from every screen. The card is rendered inside two unrelated screens, so its
scope is the union of the two, which is none of the named elements. The settings template
is rendered on its own screen and on its child screen, which is exactly
$\mathit{subtree}(\text{/settings})$. Grey arrows show component nesting and outlets.}
\label{fig:scopes}
\end{figure}

\begin{figure}[t]
\centering
\begin{tikzpicture}[x=26mm, y=15mm, every node/.style={font=\small}]
  \tikzset{el/.style={draw=muted, rounded corners=2pt, inner sep=3pt, fill=white},
           named/.style={draw=accent, fill=accentlight, rounded corners=2pt, inner sep=3pt}}
  \node[named] (bot) at (0,0) {$\varnothing = \bot$};
  \node[named] (h) at (-1,1) {$\{h\}$};
  \node[named] (s) at (0,1) {$\{s\}$};
  \node[named] (p) at (1,1) {$\{p\}$};
  \node[el, draw=warm, fill=warmlight] (hs) at (-1,2) {$\{h,s\}$};
  \node[el] (hp) at (0,2) {$\{h,p\}$};
  \node[named] (sp) at (1,2) {$\{s,p\} = \mathit{subtree}(s)$};
  \node[named] (top) at (0,3) {$\{h,s,p\} = \top$};
  \foreach \a/\b in {bot/h, bot/s, bot/p, h/hs, h/hp, s/hs, s/sp, p/hp, p/sp, hs/top, hp/top, sp/top}
    \draw[muted] (\a) -- (\b);
  \end{tikzpicture}
\caption{The scope domain $\mathcal{P}(S)$ for three screens: $h$ for \texttt{/}, $s$ for
\texttt{/settings} and its child $p$ for \texttt{/settings/profile}. Shaded elements are the
named scopes at which edges are anchored. The join is set union: a component rendered in
both $h$ and $s$ gets $\{h\} \sqcup \{s\} = \{h,s\}$, marked in orange, which is not a named element, so its edges are stored per
screen. The height of the lattice is $|S|$.}
\label{fig:lattice}
\end{figure}



\begin{table}[t]\small
\caption{Scope of each construct that can contain a navigation site.}
\label{tab:scopes}
\begin{tabularx}{\textwidth}{@{}YY@{}}
\toprule
Construct & Scope \\
\midrule
component above the router outlet (the shell) & $\top$ \\
component of route $r$ & $\mathit{subtree}(r)$ \\
component mounted at routes $r_1, \ldots, r_k$ &
$\mathit{subtree}(r_1) \sqcup \cdots \sqcup \mathit{subtree}(r_k)$ \\
component rendered in the template of another component & join of the scopes of the
components that render it. The prototype applies this row only to components that no
route mounts (Section~\ref{sec:implementation}) \\
member (method, property initialiser, lifecycle hook) of a component & the scope of the
component \\
method of an injectable service & join of the scopes of its callers \\
\texttt{canActivate}, \texttt{canActivateChild} or \texttt{canMatch} on route $r$ & join of the
scopes of the navigation sites whose target lies in $\mathit{subtree}(r)$, the screens a
navigation into $r$ starts from \\
\texttt{canDeactivate} on route $r$ & $\mathit{subtree}(r)$ \\
HTTP interceptor & $\top$ \\
NgRx effect for action $a$ & join of the scopes of every dispatcher of $a$ \\
\bottomrule
\end{tabularx}
\end{table}

The row for services is why a simple traversal of the component tree is not enough. A service has no scope of its own. Its scope is
the join of the scopes of the places that call it, which may themselves be services.
Attribution is therefore computed as a fixpoint over the input relations of
Section~\ref{sec:inputs}. The outlet relation gives the initial scopes, and scopes are
propagated along the component nesting relation $N$, from each component to its own members,
along the call graph $\mathit{CG}$, along the dispatch relation $D$, and from each navigation
site to the entry guards of the routes its target lies under, with $\mathcal{P}(S)$ as the abstract domain. This is a
monotone dataflow problem in the sense of Kildall
\cite{kildall1973unified,nielson1999principles}, and two properties of it do separate
work. Finite height gives termination. The height of $\mathcal{P}(S)$ is the number of
screens $|S|$, so the scope of one node can rise at most $|S|$ times before it reaches
$\top$, and the iteration can make only finitely many changes. Monotonicity of the
transfer functions gives the result: the iteration reaches the least solution of the
constraint system, and that solution does not depend on the order in which nodes leave
the worklist. The transfer functions here are monotone because the only operation they
apply to a scope is a join with another scope, and union is monotone in both arguments.

A worklist solver re-examines a node once for each rise of its own scope, and
the work done for a node is proportional to its out-degree in $N \cup \mathit{Mem} \cup \mathit{CG}
\cup D \cup T$, where $\mathit{Mem}$ links each component to its members and $T$ links each callable
holding a navigation site to the entry guards that navigation runs. A run therefore makes at most
\[
  O\bigl(|S| \cdot (|N| + |\mathit{Mem}| + |\mathit{CG}| + |D| + |T|)\bigr)
\]
transfer-function applications, after an initialisation that seeds each routed component with
$\mathit{subtree}(r)$ in $O(|R| \cdot |S|)$. Each application is a union of two subsets of $S$, which
costs $O(|S|)$ with scopes held as bit sets, so the running time of the solver, initialisation
included, is bounded by
$O(|R| \cdot |S| + |S|^2 \cdot (|N| + |\mathit{Mem}| + |\mathit{CG}| + |D| + |T|))$. This is the bound for the full
model, and it covers the solver only: building the relations it runs over, which includes parsing
the sources and resolving calls, is a separate cost not bounded here. $|\mathit{Mem}|$ is the number of members of components and is at most the number of callables,
and $|T|$ is at most the number of navigation sites times the depth of the route tree times
the largest number of entry guards on one route.

The part of the analysis that is implemented, the nesting, outlet, membership and call-graph
relations, is the full model with $D$ and $T$ empty, so the bound above specialises to it, and the call-graph pass, a
worklist of set-valued joins along member call edges, is bounded by it. The nesting and outlet part
is cheaper than that specialisation, because reachability from each source needs no set-valued
joins. Scopes are propagated along the component nesting
graph from the shell and from the component of each screen, which is plain reachability
and costs
\[
  O\bigl(|S| \cdot (|C| + |N|)\bigr),
\]
one traversal of the nesting graph per routed component and per shell, where $C$ is the set of
components and the number of such sources is at most $|S|$ plus the number of shells. What
makes the propagation necessary in both cases is that scopes travel transitively and
the scope of one node can rise more than once, once for each screen that reaches it.
Cycles are not the reason. The nesting, outlet and call-graph relations the prototype implements
may contain cycles, since mutually recursive templates are legal, and the worklist
handles them; a cycle in $\mathit{CG}$ or $D$ likewise only means that a node can be
re-examined after one of its own successors.

\subsection{Input relations}
\label{sec:inputs}

Four relations, in order of increasing difficulty, provide the input to the fixpoint computation. \emph{Component nesting}: $A \to B$ when $A$'s template
instantiates $B$'s selector, resolved within the correct import scope; matching
selector text globally is unsound, since two components may share a selector string in
different scopes. \emph{Router outlets}: components on the path from the root to the
outlet form the shell and take $\top$; components below the component of route $r$ take its
scope $\mathit{subtree}(r)$; a named outlet carries a second mount, which makes auxiliary routes
representable without separate treatment. \emph{Call graph with dependency injection}:
a call links to a method only after resolving the receiver's type and its binding in
the relevant injector. Resolution follows class hierarchy analysis
\cite{dean1995cha}, with Angular's provider declarations as the source of bindings, as
recommended for framework-heavy code \cite{antoniadis2020enterprise}. Where a token
admits several implementations the analysis joins over all of them, which is one
source of \emph{may} modality. \emph{Dispatch
relations}: a dispatch is treated as a call-graph edge from every dispatcher to every
handler, which admits NgRx effects and router event subscriptions under the same
fixpoint.

Algorithm~\ref{alg:attribution} gives the computation, and its propagation steps are
the transfer-function applications bounded above. A step adds a screen to a scope only
when an edge of $N$, $\mathit{Mem}$, $\mathit{CG}$, $D$ or $T$ carries it there, and scopes only grow, so
what the algorithm computes is the least solution of the constraint system that
Table~\ref{tab:scopes} states. In the model a component that is both mounted by a route
and rendered in the template of another component therefore holds the join of the two
scopes. Section~\ref{sec:implementation} reports where the prototype departs from this: it does not
propagate into a component that a route mounts, and it adds fallbacks that are not constraints
of the table, so its result is not the least solution of this system or of any restriction of it.

\begin{algorithm}[t]
\caption{Scope attribution as a worklist fixpoint over $\mathcal{P}(S)$. The relations marked \emph{model only} belong to the full analysis and are not implemented in
the prototype evaluated in Chapter~\ref{ch:results}.}
\label{alg:attribution}
\begin{algorithmic}[1]
\Require routes $R$ with route tree; for each route $r$ its entry guards $\mathit{entry}(r)$,
  named in \texttt{canActivate}, \texttt{canActivateChild} or \texttt{canMatch}, and its
  \texttt{canDeactivate} guards $\mathit{leave}(r)$; components $C$; nesting
  $N \subseteq C \times C$; HTTP interceptors $I$; call graph $\mathit{CG}$; dispatch relation
  $D$; guard relation $T$, linking a callable $x$ to $g \in \mathit{entry}(r)$ when $x$ holds a
  navigation site whose target lies in $\mathit{subtree}(r)$
\Ensure $\mathit{scope}(x) \in \mathcal{P}(S)$ for every component and callable $x$; the
  callables include the members of components and of services
\ForAll{$x$} $\mathit{scope}(x) \gets \varnothing$ \EndFor
\ForAll{$r \in R$ with $\mathit{component}(r) = c$}
  \State $\mathit{scope}(c) \gets \mathit{scope}(c) \cup \mathit{subtree}(r)$
  \Comment{own screen and outlet descendants}
\EndFor
\ForAll{$c$ on the nesting path from the root component to the holder of the root-level primary outlet}
  \State $\mathit{scope}(c) \gets S$ \Comment{application shell, $\top$}
\EndFor
\ForAll{$r \in R$ and $g \in \mathit{leave}(r)$}
  \State $\mathit{scope}(g) \gets \mathit{scope}(g) \cup \mathit{subtree}(r)$
  \Comment{runs as the user leaves $r$}
\EndFor
\ForAll{HTTP interceptor $i \in I$} \Comment{\emph{model only}: registered for the application}
  \State $\mathit{scope}(i) \gets S$
\EndFor
\State $W \gets$ every $x$ with $\mathit{scope}(x) \neq \varnothing$
\While{$W \neq \emptyset$}
  \State take $x$ from $W$
  \ForAll{$y$ with $(x,y) \in N$} \Comment{$x$'s template renders $y$}
    \State \Call{Propagate}{$x, y$}
  \EndFor
  \ForAll{members $m$ of $x$, if $x$ is a component} \Comment{a member runs where its component is shown}
    \State \Call{Propagate}{$x, m$}
  \EndFor
  \ForAll{$y$ with $(x,y) \in \mathit{CG} \cup D \cup T$} \Comment{$\mathit{CG}$: member level, bindings read application-wide; $D$, $T$: \emph{model only}}
    \State \Call{Propagate}{$x, y$}
  \EndFor
\EndWhile
\Procedure{Propagate}{$x, y$}
  \If{$\mathit{scope}(x) \not\subseteq \mathit{scope}(y)$}
    \State $\mathit{scope}(y) \gets \mathit{scope}(y) \cup \mathit{scope}(x)$; add $y$ to $W$
  \EndIf
\EndProcedure
\end{algorithmic}
\end{algorithm}


\subsection{Guards as program points}

A guard that returns a \texttt{UrlTree} is itself a navigation, typically to a login screen.
Its source is not the guarded route: a \texttt{canActivate} or \texttt{canMatch} guard on $r$
runs while the router leaves the screen the navigation started from, so the redirect edge
leaves that screen, the origin of an edge into $\mathit{subtree}(r)$, and only a
\texttt{canDeactivate} guard redirects from a screen of $\mathit{subtree}(r)$ itself. Treating guards only as labels on
inbound edges misses this edge class entirely. Further, the guard of an edge in the model is the
conjunction of the conditions on the containment path from screen to call site
(structural directives) with the guards the router runs on that navigation
(Section~\ref{sec:model}). Reporting only the second part gives a weaker condition than the true
one. The prototype does not materialise that conjunction in the graph that is measured. There,
a structural condition in a template, or an enclosing \texttt{if} or conditional expression
around a navigation call, demotes the edge to may but is not written into its guard, so the guard
of a screen-to-screen edge holds only the router's guards. Only the optional mode that anchors
navigations at the controls that trigger them, which no table of Chapter~\ref{ch:results} uses,
keeps the enclosing condition of a handler's navigation as source text beside the router's guards.
The graph evaluated in Chapter~\ref{ch:results} therefore records that a navigation is
conditional without recording the condition. In the baseline, for example,
a link that is rendered only inside \texttt{@if (isUser())} was labelled with the
authentication guard of its target route. That guard is satisfied by any logged-in
user, including users who are not the owner and who therefore never see the link.

\subsection{Boundary}
\label{sec:boundary}

A navigation inside a signal effect has a firing condition that is a data dependence
on a reactive graph. As Section~\ref{sec:angular} notes, only the read set is
statically visible. The model gives such edges the modality \emph{unknown} and records
the set of signals read by the effect as the symbolic guard; no pass in this version
emits one. The analysis does not try to determine when these edges fire. It marks this
limit explicitly instead of giving a result that looks more precise than it is.

\subsection{Anchored edges}
\label{sec:anchored}

If the result of attribution were stored explicitly, a single link in the application
shell would produce one edge for every screen. Where every site has a single target, the size of
the graph would then be bounded by $|S| \cdot |\mathit{sites}|$, where $\mathit{sites}$ is the
set of navigation sites; for the largest application in
the corpus that bound is $\RoundTwoEarlierOracleMifosNow{} \times \ConstructsMifos{}$. A site
whose target is only partly known can reach several screens, so in general
\[
  |E_{\mathit{expanded}}| \;\le\; \sum_{x \in \mathit{sites}} |\mathit{scope}(x)| \cdot |\mathit{Targets}(x)|,
\]
which is $O(|S|^2 \cdot |\mathit{sites}|)$ when a site of scope $\top$ reaches every screen.
Anchoring removes the factor that comes from the source scope, not the one that comes from the
targets. Instead, each edge is stored once, anchored
at its scope, and the set of edges leaving a particular screen is computed when it is
requested:
\[
  \mathit{edgesFrom}(r) \;=\; E_\top \;\cup\; \bigcup_{r' \succeq r}
  E_{\mathit{subtree}(r')} \;\cup\; E_{\{\mathit{screen}(r)\}},
\]
where $E_X$ is the set of edges anchored at scope $X$ and $r' \succeq r$ is the
reflexive ancestor relation on routes introduced above, so that the term for
$r' = r$ contributes $E_{\mathit{subtree}(r)}$ itself. The query is made for a screen,
so $\mathit{component}(r)$ is defined, while $r'$ ranges over all ancestors of $r$,
prefix routes included. An edge whose scope is some other
screen set is stored at each of its screens and reaches $\mathit{edgesFrom}(r)$ through
the $E_{\{\mathit{screen}(r)\}}$ term.
Figure~\ref{fig:anchored} contrasts the two representations.

\begin{figure}[t]
\centering
\begin{tikzpicture}[node distance=5mm]
  % materialised
  \node[font=\small\itshape] at (0,2.4) {materialised};
  \node[screen, minimum width=14mm] (a1) at (-1.6,1.4) {/};
  \node[screen, minimum width=14mm] (a2) at (-1.6,0.6) {/feed};
  \node[screen, minimum width=14mm] (a3) at (-1.6,-0.2) {/editor};
  \node[screen, minimum width=14mm] (a4) at (-1.6,-1.0) {/profile};
  \node[screen, minimum width=16mm] (at) at (1.6,0.2) {/settings};
  \foreach \n in {a1,a2,a3,a4} \draw[mayedge] (\n.east) -- (at.west);
  \node[font=\scriptsize, text=muted] at (0,-1.7) {one edge per screen, $|S|$ in all};

  % anchored
  \begin{scope}[xshift=70mm]
  \node[font=\small\itshape] at (0,2.4) {anchored};
  \node[screen, minimum width=14mm] (b1) at (-1.6,1.4) {/};
  \node[screen, minimum width=14mm] (b2) at (-1.6,0.6) {/feed};
  \node[screen, minimum width=14mm] (b3) at (-1.6,-0.2) {/editor};
  \node[screen, minimum width=14mm] (b4) at (-1.6,-1.0) {/profile};
  \node[draw=accent, fill=accentlight, rounded corners=6pt, font=\small, inner sep=3pt] (top) at (0,0.2) {$\top$};
  \node[screen, minimum width=16mm] (bt) at (1.8,0.2) {/settings};
  \foreach \n in {b1,b2,b3,b4} \draw[accent, densely dotted] (\n.east) -- (top);
  \draw[mayedge] (top) -- (bt);
  \node[font=\scriptsize, text=muted] at (0,-1.7) {one edge; members expanded on query};
  \end{scope}
\end{tikzpicture}
\caption{The same header link, rendered under a condition and so held at may on every screen,
stored two ways. Materialised, it becomes one edge per screen, including one from
\texttt{/settings} itself, which is not drawn. Anchored, it is stored once at the scope node $\top$,
and $\mathit{edgesFrom}(r)$ adds it to the edges of any screen $r$ when asked.
Claim~\ref{claim:expansion} states that both forms describe the same transitions. Only a reason
for may that every screen shares keeps the anchored form: an unconditional link is demoted from
\texttt{/settings}, the screen it already displays, but not from the others, so the extractor
stores it as one edge per screen.}
\label{fig:anchored}
\end{figure}


\begin{claim}[Expansion]
\label{claim:expansion}
The expansion of the anchored edge set equals the materialised edge set.
\end{claim}

\noindent The claim is tested rather than proved: Section~\ref{sec:res-after} compares the two
sets on the corpus and on fixtures.

\section{Implementation}
\label{sec:implementation}

The extractor, \emph{uigraph}, is available as open source. It consists of a core
that is independent of any particular framework and a set of adapters, one per
framework. The core defines the graph model, its representation on disk, and the
operations over graphs. Each adapter maps the source constructs of one framework onto
the shared model, following a fixed interface. The Angular adapter reads route
configuration, templates and class code with the TypeScript compiler API, and emits
nodes for screens and edges for \texttt{routerLink} attributes and
\texttt{Router.navigate} calls, each with the rule identifier and source location that
produced it as witness.

Scope attribution is implemented as a separate pass in the Angular adapter, enabled
with the option \texttt{--attribution}. With the option off, the extractor produces
exactly the baseline edge set; this was checked on every application in the corpus.
The pass implements three of the four input relations of Section~\ref{sec:attribution}:
component nesting, router outlets and, as the second departure below describes, the call graph.
The dispatch relation is not implemented.

\textbf{Component nesting.} Every component in the application, excluding test hosts
and stories, is read with its element selectors and its template. An edge $A \to B$
is added when $A$'s template contains an element with $B$'s selector. When a selector
belongs to more than one component, the candidates are narrowed to those that $A$'s
file imports; if none is imported, the analysis joins over all of them and records a
note. Edges from such a join are not capped at \emph{may}, so a site in one of the candidates
can be placed on a screen that renders another; Claim~\ref{claim:must} excludes such edges
(Section~\ref{sec:approximations}); whether any corpus must edge comes from one was not measured. An attribute selector is matched the same way: a directive is rendered by the
templates that carry its attribute, so it has a scope in the same way as a component does.

\textbf{Router outlets.} The route tree is recorded during route extraction. A route
component's template is rendered on its own screen and on every descendant screen, so
its navigation sites get the scope $\mathit{subtree}(r)$. A component that is not
routed, is not placed on a screen by the nesting relation and holds the primary
\texttt{<router-outlet>} is the application shell, and everything it renders gets the
scope $\top$. A component that is not routed and not referenced, but whose subtree
contains the outlet, is most likely mounted by a route the extractor cannot see, for
example through a helper function that wraps a route table. It is also treated as
shell, but its edges are capped at \emph{may}, so that Claim~\ref{claim:must} is not
put at risk by a scope that may be too large.

The pass departs from Algorithm~\ref{alg:attribution} in two respects. The first is that it does not
propagate a scope into a component that a route mounts: a routed component keeps
$\mathit{subtree}(r)$ and does not also receive the scope of a component whose template
renders it. Traversal does pass through a routed component, so a component nested beneath it still receives
the scopes of the components above; only the routed component's own sites are left at
$\mathit{subtree}(r)$. For the nesting and outlet relations the result is contained in the
solution the full system would give, so a navigation site in a routed component can be
attributed to fewer screens than the model allows, never to more. A scope that is too small loses edges and cannot invent them, which is why the
restriction cannot put Claim~\ref{claim:must} at risk; what it costs is recall, and the
constructs it strands are counted in Chapter~\ref{ch:results}.

The second departure is in how the call graph of Section~\ref{sec:inputs} is read.
Algorithm~\ref{alg:attribution} ranges over every component and callable, and the pass does too:
its nodes are the members of classes, and a member of a component runs on that component's
screens. A call \texttt{this.f.g()}, where \texttt{f} is a
constructor parameter, a typed field, a field initialised with \texttt{inject(T)}, or an
\texttt{@Inject(TOKEN)} parameter, is resolved through the declared type of \texttt{f} to the class that type names, to the classes a provider
declaration binds it to with \texttt{useClass} or \texttt{useExisting}, or, for an abstract class
no provider binds, to its subclasses in the application, which is class hierarchy analysis
\cite{dean1995cha} restricted to the declared type. The departure is that bindings are read across
the whole application rather than relative to the injector that resolves them, so where two
injectors bind a token differently the join covers both; that can only widen a scope. A factory or value provider is not read: a call on its token resolves to the token's own
class if the application declares one, and a call on a type declared outside the application
contributes no edge. The seeded call-graph iteration on its own reaches the least solution of
its constraints. The scopes the extractor emits for service sites add two fallbacks to that
solution. A member no caller reaches falls back to the join over
every component that injects its class, the per-class approximation an earlier version used for
every service. A member the iteration does reach can still receive more screens: where a caller
the graph does not see, one that nothing calls and nothing seeds, reaches the member, the join of
that caller's class is added, or this class's own join if nothing places the caller. Neither
fallback is a constraint of Table~\ref{tab:scopes}, and the first is not
monotone in the call graph: a member whose only caller is found receives that caller's screens
instead of the wider join, so adding a call edge can narrow its scope. The emitted service
scopes are therefore the least call-graph solution widened by these fallbacks, for reached
members as well as for unreached ones, and not the least solution of the constraints in Table~\ref{tab:scopes}. Service sites are held at \emph{may} either way,
because the calls on the way carry conditions the analysis does not follow.

\textbf{Anchored edges.} Edges with scope $\top$ are stored once at a node of kind
\texttt{scope}. Edges whose screen set is exactly $\mathit{subtree}(r)$ are stored once
at a scope node that lists the members of that subtree. Other screen sets are stored
as one edge per screen. The core library provides the expansion query
$\mathit{edgesFrom}(r)$ and a function that materialises all anchored edges, which is
used to test Claim~\ref{claim:expansion}.

The dispatch relation is not implemented in this version, so a navigation in an NgRx effect has
no scope from the model itself. Two approximations stand in where the model has no relation the
pass implements, and Section~\ref{sec:res-ablation} withholds them below the \emph{full} level, so
no lower level counts their sites. The relation $T$ is not implemented, and a guard is given
instead the screens of the routes that name it, as a bare identifier, in \texttt{canActivate},
\texttt{canActivateChild}, \texttt{canMatch} or \texttt{canDeactivate}: the screen of each such
route, not $\mathit{subtree}(r)$, and nothing for a route that is no screen, such as a
componentless prefix. For a \texttt{canDeactivate} guard that is narrower than the model's
scope; for an entry guard it is not the model's rule, since the redirect leaves the screen a navigation starts from, so the
stand-in can be too small where navigations into the route start elsewhere and too large where
they do not start at the route itself. An effect, being registered for the lifetime of the
application, is given $\top$. Edges from both are capped at \emph{may}. A class the call graph
does not reach and nothing injects still has no scope, and the evaluation counts those in the
defect class ``Attribution, service'' of Table~\ref{tab:loss-rules}. The pass is
covered by unit tests for each of the cases above, including a test that the baseline
is unchanged when the option is off.

One change was made to the baseline itself. For components with an inline template,
the witness of an edge pointed at the line where the template starts rather than at
the construct. The witness now points at the construct's own line. This can change
witness locations but no edge. On the corpus, the baseline graphs with and without this
change are identical, including witnesses.
